\documentclass[
  aps,
  prb,
  reprint,
  amsmath,
  amssymb,
  superscriptaddress,
  longbibliography,
  nofootinbib,
  floatfix
]{revtex4-2}

\usepackage[T1]{fontenc}
\usepackage{lmodern}
\usepackage{graphicx}
\usepackage{bm}
\usepackage{microtype}
\usepackage[hidelinks]{hyperref}

\graphicspath{{figures/}}

\newcommand{\Heff}{H_{\mathrm{eff}}}
\newcommand{\Jeff}{J_{\mathrm{eff}}}
\newcommand{\ceff}{c_{\mathrm{eff}}}
\newcommand{\qpeak}{q_{\mathrm{peak}}}
\newcommand{\Szz}{S^{zz}(q,\omega)}
\newcommand{\avg}[1]{\left\langle #1\right\rangle}
\newcommand{\pentamer}{pentamer}

\begin{document}
\raggedbottom

\title{Local-spin projection and momentum-selective spectral weight
in a Heisenberg pentamer chain}

\author{Author Name}
\email{corresponding.author@example.edu}
\affiliation{Department and Institution, City, Postal Code, Country}

\date{\today}

\begin{abstract}
Different physical momenta can probe the same collective mode in a chain of
spin clusters, yet carry very different spectral weights.
We study this distinction in a spin-$1/2$ Heisenberg chain with a five-site unit
cell and reflection-symmetric intra-cell exchange.
Changing the outer intra-cell bonds from antiferromagnetic to ferromagnetic
changes the isolated-cell ground multiplet from spin $1/2$ to spin $3/2$ in
the sampled range. The projected spin distribution then reverses the relative
brightness of the folded antiferromagnetic responses at $q=0.60\pi$ and $\pi$.
For two weakly coupled cells we derive the absolute low-energy spectral weight
without fitting an amplitude. Exact diagonalization gives a relative $L_2$
error below $0.27\%$ for the momentum-resolved weight vector and verifies
the projected singlet--triplet gap to within
$0.56\%$ at the six tested points with
$J_3/\Delta_{\mathrm{loc}}\le0.01$.
Reflection parity independently makes the lowest excited local doublet dark
at $q=\pi$ for positive intra-cell exchange.
Complete $L=40$ spectra at stronger coupling compare the redistribution
of spectral weight at a common finite-time resolution, outside the quantitative
domain of the weak-coupling projection.
Static data at $L=60,80$ resolve a momentum reconstruction on the ferromagnetic
side, but do not identify a common phase-boundary interval.
The results distinguish changes in the visibility of collective spin
excitations from evidence for a new thermodynamic phase.
\end{abstract}

\maketitle

\section{Introduction}
\label{sec:introduction}

A spin cluster carries more information than its total spin. The distribution
of that spin over the sites determines how neighboring clusters interact and
how a momentum-resolved probe couples to their collective motion. In a chain
of clusters, several physical momenta can describe the same effective-chain
momentum, yet their scattering intensities need not be comparable. This makes
spectral weight, rather than the excitation energy alone, a useful probe of
the internal structure of a composite spin.

Trimer chains provide an established example: their ground-state doublets form
an effective spin-$1/2$ chain, while perturbative and numerical spectra connect
higher-energy features to internal trimer excitations~\cite{cheng2022trimer}.
Studies of tetramer chains similarly combine cluster states, effective
descriptions, and dynamical structure factors to distinguish collective and
local excitations~\cite{li2026tetramer}.
The relation between cluster symmetry and the visibility of a transition is
also known for isolated pentamers~\cite{haraldsen2011pentamer}.
These results motivate a specific question for a chain whose cluster spin can
change: how does that change redistribute low-energy intensity among the
physical momenta representing the same collective fluctuation?

We address this question in a linear spin-$1/2$ Heisenberg chain with the
five-bond pattern $J_2$--$J_1$--$J_1$--$J_2$--$J_3$.
The first four bonds define a reflection-symmetric pentamer, and $J_3>0$
couples neighboring cells. We set $J_1=1$ and write
$\alpha=J_2/J_1$ and $\beta=J_3/J_1$.
In the isolated-cell ranges studied, positive $\alpha$ gives a spin-$1/2$
ground doublet, whereas negative $\alpha$ gives a spin-$3/2$ quartet.
Both are half-integer degrees of freedom; the distinction is not the
integer--half-integer contrast between an even and an odd cell.
The Lieb--Schultz--Mattis constraint rules out a unique, symmetry-preserving
gapped bulk ground state for a half-integer spin per translation cell under
its stated symmetry assumptions~\cite{lieb1961,ogata2019}, but it does not fix
the momentum dependence of the response or establish a particular critical
theory throughout this model's coupling plane.

Projecting the five site spins onto the cell ground multiplet gives
$P\bm S_aP=c_a\bm T$. The coefficients $c_a$ determine both the leading
exchange $\Jeff=\beta c_1c_5$ and the interference factor
$F(q)=|\sum_a c_a e^{iqa}|^2$ seen by the probe.
Because the cell momentum is $Q=5q\pmod{2\pi}$, the physical momenta
$q=0.2\pi$, $0.6\pi$, and $\pi$, together with their reflection partners,
all represent $Q=\pi$.
For $\alpha=0.5$ the cell factor strongly favors $q=\pi$, whereas for
$\alpha=-0.5$ it favors $q=0.6\pi$ and nearly cancels at $q=\pi$.
We derive the absolute low-energy spectral weight of two coupled cells and
test it by exact diagonalization over a fixed weak-coupling sequence.
This establishes a quantitative relation between the cell's internal spin
distribution and the relative visibility of folded responses, with an
explicitly tested range of validity.

Reflection provides a second, distinct constraint on the spectrum.
The isolated positive-$\alpha$ cell has an odd-parity lowest excited doublet
that an even $q=\pi$ spin operator cannot reach. Thus the first excited local
level need not be the first visible one.
We compare these local predictions with ground-state and dynamical
matrix-product-state calculations~\cite{white1992,schollwoeck2011,paeckel2019}.
The positive-coupling static data support an effective antiferromagnetic chain
at the tested points. Negative couplings produce a finite-size reconstruction
of the dominant static momentum, and complete connected $L=40$ spectra at
$(\alpha,\beta)=(\pm0.5,0.5)$ compare the two responses at a common
finite-time resolution. These latter points lie outside the tested
weak-coupling regime; they show how spectral weight is organized at finite
coupling, rather than verifying the projected weights there or establishing
a new thermodynamic phase.

Section~\ref{sec:methods} defines the model and numerical observables.
Section~\ref{sec:results} derives and tests the projected spectral weights;
Sec.~\ref{sec:dynamics} presents the finite-coupling chain results.
Section~\ref{sec:discussion} explains the relation between the two local-spin
descriptions and the reconstruction near $\alpha=0$, and identifies the
spectroscopic consequences that follow within the controlled limits.

\section{Model and observables}
\label{sec:methods}

\subsection{Microscopic chain and cell momentum}

The chain contains $N$ complete five-site cells and $L=5N$ spin-$1/2$ sites,
with Hamiltonian
\begin{align}
H={}&J_1\sum_n(\bm S_{n,2}\!\cdot\!\bm S_{n,3}
 +\bm S_{n,3}\!\cdot\!\bm S_{n,4})\notag\\
&+J_2\sum_n(\bm S_{n,1}\!\cdot\!\bm S_{n,2}
 +\bm S_{n,4}\!\cdot\!\bm S_{n,5})\notag\\
&+J_3\sum_n\bm S_{n,5}\!\cdot\!\bm S_{n+1,1}.
\label{eq:hamiltonian}
\end{align}
We set $J_1=1$, $\hbar=1$, and physical-site spacing to one, and write
$\alpha=J_2/J_1$ and $\beta=J_3/J_1>0$.
Open chains terminate at complete-cell boundaries.
The isolated-cell Hamiltonian is invariant under $a\mapsto6-a$.
Its 32 states are resolved by total spin and reflection parity;
transition weights are summed over complete multiplets, so they do not depend
on the eigenvector basis chosen within a degenerate level.

The distinction between physical momentum $q$ and cell momentum
$Q=5q\pmod{2\pi}$ is central to this work.
For example, $q=0.20\pi$, $0.60\pi$, and $\pi$ all correspond to $Q=\pi$,
but their phases within a cell differ.
They therefore need not have equal intensity even when they probe the same
effective-chain excitation.
Figure~\ref{fig:model} illustrates the cell and its positive-$\alpha$ projection.

\begin{figure*}[t]
  \centering
  \includegraphics[width=0.94\textwidth]{fig1_model_projection.png}
  \caption{Pentamer chain and isolated-cell structure.
  The outer cell bonds have exchange $J_2$, the central bonds $J_1$, and
  neighboring cells are connected by $J_3$.
  The local levels and spin-$1/2$ projection shown here refer to positive
  $\alpha=J_2/J_1$; the spin-$3/2$ projection for negative $\alpha$ is
  developed in the text. All energies are in units of $J_1$.
  Plotted levels come from numerical
  diagonalization.}
  \label{fig:model}
\end{figure*}

\subsection{Static and dynamical probes}
\label{sec:diagnostics}

We use distinct static and dynamical observables, with their boundary
conventions stated explicitly.
The static scans employ the connected central-window structure factor
\begin{equation}
 S_W^{zz}(q)=\frac{1}{M}\sum_{i,j\in W}e^{iq(j-i)}
 [\langle S_i^zS_j^z\rangle-\langle S_i^z\rangle\langle S_j^z\rangle],
 \label{eq:static}
\end{equation}
where $W$ removes $\lfloor L/10\rfloor$ sites from each end and $M=|W|$.
For negative $\alpha$, $\qpeak$ maximizes this quantity on 401 momenta in
$[0,\pi]$.
Its local sharpness is
\begin{equation}
 R_L=1-\frac{S_W^{zz}(\qpeak-\delta)+S_W^{zz}(\qpeak+\delta)}
 {2S_W^{zz}(\qpeak)},\quad \delta=\frac{2\pi}{M},
 \label{eq:ratio}
\end{equation}
using the nearest sampled neighbors. At $0$ or $\pi$ the ratio is one-sided.
The size flow $\Delta R=R_{80}-R_{60}$ is an operational measure of peak
sharpening, not by itself a phase invariant.

The full-momentum dynamical maps instead use every site, with no spatial window:
\begin{equation}
 O_q=\frac{1}{\sqrt L}\sum_{j=1}^Le^{iqj}
 (S_j^z-\langle0|S_j^z|0\rangle),
 \label{eq:momentum-source}
\end{equation}
\begin{equation}
 S^{zz}(q,\omega)=\sum_m|\langle m|O_q|0\rangle|^2
 \delta(\omega-E_m+E_0).
 \label{eq:dsf}
\end{equation}
This normalization is also used in the two-cell exact-diagonalization test.
It fixes the absolute spectral weight rather than just its momentum dependence.

\subsection{Finite-chain calculations}
\label{sec:convergence}
\label{sec:dsf}

Ground states are obtained by density-matrix renormalization group (DMRG) in
fixed total-$S^z$ sectors, using matrix-product states
(MPS)~\cite{white1992,schollwoeck2011}.
The negative-$\alpha$ static scans use two distinct $S^z=0$ initial states and a
separate $S^z=1$ calculation. Retained runs satisfy variance per site
$\mathrm{Var}(H)/L\le5\times10^{-8}$, final truncation error
$\le10^{-9}$, sector magnetization error $\le10^{-8}$, and
$E_0(S^z=1)-E_0(S^z=0)\ge-10^{-8}$, with at most 64 sweeps.
These are numerical screening criteria, not rigorous error bounds.
Appendix~\ref{sec:krylov} examines a case where passing them did not ensure
stability of an entropy fit against the local eigensolver setting.

For the dynamics, real-time MPS evolution~\cite{vidal2004,paeckel2019} gives
$C_q(t)=\langle b_q|e^{-i(H-E_0)t}|b_q\rangle$ for
$|b_q\rangle=O_q|0\rangle$.
The displayed estimator is
\begin{equation}
 S_T^{zz}(q,\omega)=\frac{1}{\pi}\operatorname{Re}
 \int_0^Tdt\,e^{i\omega t}(1-t/T)C_q(t).
 \label{eq:finite-time-dsf}
\end{equation}
The triangular time window, including the actual $t=0$ value, is part of the
estimator. Its line kernel has full width at half maximum
$5.56623/T$; sampling $\omega$ more finely does not improve this resolution.
The two $L=40$ maps use $TJ_1=60$, second-order odd--even time evolution with
$\Delta tJ_1=0.05$, maximum MPS bond dimension $\chi=192$, cutoff $10^{-10}$,
and 41 independently evolved sources $q_m=2\pi m/40$.
No Gaussian smoothing, momentum interpolation, reflection averaging, or
spectral clipping is applied.
The validation and display normalization are specified in
Appendix~\ref{sec:spectral-control}.

The same deterministic calculation routines supply the figures and the
source-data tables. Exact diagonalization, projection identities, source-norm checks, and the
parameter controls below are the recorded validation procedures.

\section{Local spin and momentum-dependent visibility}
\label{sec:results}
\label{sec:doublet}

\subsection{Projection onto a cell spin}

The isolated cell retains a half-integer spin for either sign of the outer
bonds. All 100 sampled points with $0<\alpha\le1$ have an even-parity
spin-$1/2$ ground doublet; all 201 sampled points with
$-1\le\alpha\le-0.001$ have a spin-$3/2$ ground quartet.
The different local spins provide two realizations of an antiferromagnetically
coupled cluster chain, with different internal spin distributions.

Let $P$ project onto one isolated ground multiplet of spin $s$.
Within this subspace rotational symmetry gives
\begin{equation}
 P\bm S_{n,a}P=c_a\bm T_n,\qquad \sum_{a=1}^5c_a=1,
 \label{eq:projection}
\end{equation}
where $\bm T_n$ is a spin-$s$ operator.
The coefficients $c_a$ specify where the retained spin resides within a cell;
negative coefficients indicate a local moment antiparallel to the cell spin.
Reflection implies $c_a=c_{6-a}$ for either reflection parity, because the
parity phases of the bra and ket cancel.
Consequently the first-order Hamiltonian is
\begin{align}
 \Heff&=NE_5+\Jeff\sum_n\bm T_n\!\cdot\!\bm T_{n+1},\notag\\
 \Jeff&=\beta c_1c_5=\beta c_1^2>0.
 \label{eq:effective}
\end{align}
Here $E_5$ is the cell ground-state energy; corrections are of order
$\beta^2/\Delta_{\mathrm{loc}}$ at fixed cell parameters, where
$\Delta_{\mathrm{loc}}$ is the gap above the entire ground multiplet.
The controlled limit is $g=\beta/\Delta_{\mathrm{loc}}\ll1$.

Projecting the probe, rather than only the Hamiltonian, reveals the observable
consequence. Its intra-cell amplitude is
$f(q)=\sum_ac_ae^{iqa}$ and its shape factor is
\begin{equation}
 F(q)=|f(q)|^2.
 \label{eq:shape}
\end{equation}
With the physical and effective probes normalized by $1/\sqrt L$ and
$1/\sqrt N$, respectively,
\begin{equation}
 S^{zz}_{\mathrm{phys}}(q,\omega)\simeq
 \frac{F(q)}5 S^{zz}_{\mathrm{eff}}(Q,\omega),
 \quad Q=5q\pmod{2\pi}.
 \label{eq:folded-response}
\end{equation}
The factor $1/5$ is fixed by normalization, not fitted.
Thus replicas at the same $Q$ have a common collective response but unequal
visibility set by $F(q)$.

Table~\ref{tab:projection} makes this contrast quantitative at
$\alpha=\pm0.5$.
For positive $\alpha$ the $q=\pi$ replica is brightest; for negative $\alpha$
it is suppressed by almost three orders of magnitude relative to $0.60\pi$.
The latter is the largest of the three folded antiferromagnetic replicas,
not the global maximum of $F(q)$: $F(0)=(\sum_ac_a)^2=1$.
The location of a peak in the full spectrum depends on the product in
Eq.~\eqref{eq:folded-response}, not on $F$ alone.

\begin{table}[!htbp]
 \caption{Isolated-cell projection at the two representative couplings.
 Reflection fixes $c_5=c_1$ and $c_4=c_2$. The last three rows compare
 physical momenta with the same cell momentum $Q=\pi$.}
 \label{tab:projection}
 \begin{ruledtabular}
 \begin{tabular}{lrr}
 & $\alpha=0.5$ & $\alpha=-0.5$\\
 $s$ & $1/2$ & $3/2$\\
 $\Delta_{\mathrm{loc}}/J_1$ & 0.342371 & 0.126397\\
 $c_1$ & 0.622839 & 0.331329\\
 $c_2$ & $-0.277350$ & 0.242242\\
 $c_3$ & 0.309022 & $-0.147142$\\
 $\Jeff/\beta$ & 0.387928 & 0.109779\\
 $F(0.20\pi)$ & 0.060121 & 0.202128\\
 $F(0.60\pi)$ & 0.278088 & 0.693819\\
 $F(\pi)$ & 4.449570 & 0.000963\\
 \end{tabular}
 \end{ruledtabular}
\end{table}

\subsection{An absolute two-cell spectral-weight test}
\label{sec:projection-benchmark}

A two-cell chain provides a direct test of both the projected energy scale
and the probe normalization, without assumptions about a thermodynamic line
shape. Two antiferromagnetically coupled spins $s$ form a singlet, with
$\langle(T_1^z)^2\rangle=s(s+1)/3$ and
$\langle T_1^zT_2^z\rangle=-s(s+1)/3$.
Their projected physical probe is
$f(q)[T_1^z+e^{i5q}T_2^z]/\sqrt{10}$.
Its norm and the singlet--triplet splitting therefore predict
\begin{align}
 W_{\mathrm{low}}^{(1)}(q)
 &=\frac{F(q)s(s+1)}{15}[1-\cos(5q)],
 \label{eq:two-cell-weight}\\
 \Delta_{10}^{(1)}&=\beta c_1c_5.
 \label{eq:two-cell-gap}
\end{align}
There is no adjustable amplitude or independent normalization at each momentum.

We compare these expressions with complete exact diagonalization in the
$S^z=0$ sector (dimension 252) of two full pentamers and the lowest energy in
$S^z=1$ (dimension 210).
Before diagonalization, the low band is defined as
$0<\omega<\Delta_{\mathrm{loc}}/2$.
The calculation uses $g=0.0025,0.005,0.01,0.02,0.04$ at each sign of $\alpha$;
$g\le0.01$ is the predefined weak domain.
Both the spectral completeness sum rule and the Fourier-averaged local moment
are checked to absolute error below $10^{-11}$.

Figure~\ref{fig:projection-benchmark} shows that Eq.~\eqref{eq:two-cell-weight}
reproduces the dominant momentum pattern and its suppressed replicas.
Across the six weak-domain cases the largest gap error is
$0.5574\%$ and the largest relative $L_2$ error of the low-weight vector is
$0.2643\%$, below the predefined $2\%$ and $3\%$ tolerances.
At $g=0.01$, the positive and negative weight errors are $0.2643\%$ and
$0.0884\%$, respectively.
The vector error is not a uniform pointwise bound: relative errors can be
larger where destructive interference makes the predicted intensity tiny.
This test verifies absolute low-band weights as well as the effective energy
scale; the reversal of replica brightness is not a consequence of independently
rescaling the spectra.

\begin{figure*}[t]
 \centering
 \includegraphics[width=0.96\textwidth]{fig_projection_ed_benchmark.pdf}
 \caption{Absolute low-energy projection test for two pentamers ($L=10$).
 (a,b) Low-band connected weights at $g=0.01$ for $\alpha=0.5$ and $-0.5$.
 Circles are exact-diagonalization values at 41 momenta; lines join
 Eq.~\eqref{eq:two-cell-weight} at the same momenta, without a fitted amplitude.
 The band $0<\omega<\Delta_{\mathrm{loc}}/2$ is fixed before the calculation.
 (c,d) Gap error $|\Delta_{\mathrm{ED}}/\Delta_{10}^{(1)}-1|$ and weight error
 $\|\mathbf W_{\mathrm{ED}}-\mathbf W^{(1)}\|_2/\|\mathbf W^{(1)}\|_2$.
 Shading marks the predefined domain $g\le0.01$; segments connect only the
 five sampled couplings. This small-system test does not establish
 thermodynamic convergence. All symbols
 and lines are rendered from the calculated data and analytical prediction.}
 \label{fig:projection-benchmark}
\end{figure*}

\subsection{A distinct parity filter for local excitations}
\label{sec:parity}

The low-band interference above and the selection of excited local multiplets
are different effects.
At $q=\pi$, the alternating signs within a five-site cell are invariant under
$a\mapsto6-a$. The local spin probe is therefore reflection-even and cannot
connect states of opposite parity.
For positive $\alpha$ the ground doublet is even and the lowest excited
doublet is odd. This excited doublet is dark, although it is the lowest local
excitation in energy.
At $\alpha=0.14,0.23,0.33$, its multiplet-summed weight is below $10^{-25}$,
whereas spin-$3/2$ multiplets carry $84\%$--$98\%$ of the local inelastic
weight (Appendix~\ref{sec:local-details}).
Reflection selects which local states can be reached; it does not alone
identify a dispersive band after the cells are strongly coupled.

\section{Finite-coupling spectra and momentum reconstruction}
\label{sec:dynamics}

\subsection{Full-momentum response at fixed resolution}

Figure~\ref{fig:fullq} compares the connected spectra at
$(\alpha,\beta)=(0.5,0.5)$ and $(-0.5,0.5)$ for $L=40$.
At positive $\alpha$, low-energy weight is prominent near $q=\pi$ and coexists
with higher-energy bands. At negative $\alpha$, the low-energy response shifts
toward other momenta and the higher-energy pattern also changes.
All 41 momentum sources are independently calculated, and their reflection
deviations are only $2.87\times10^{-6}$ and $4.49\times10^{-6}$ in relative
root-mean-square norm. The symmetry is thus a numerical check, not imposed
by averaging the left and right halves.

These spectra extend the comparison to substantial inter-cell coupling, but
do not extend the quantitative validity of the weak-coupling formula.
Here $g=1.4604$ and $3.9558$ for the two signs.
Two-cell diagonalization at the same $\beta=0.5$ fails the predefined
low-band separation criterion for both signs; the first-order gap errors
increase to approximately $10\%$ and $84\%$.
The weak-limit replica weights therefore provide a reference for reading the
redistribution, not a quantitative fit or a unique assignment of these bands.

The energy resolution is $0.09277J_1$.
Independent sector calculations give open-chain gaps $0.05562J_1$ and
$0.001092J_1$, both narrower than this width.
The maps can compare weight distributions on the displayed energy scale,
but do not establish that either gap is resolved from zero frequency.
Weak stripes enhanced by the common color scale must likewise not be counted
as additional excitations on their appearance alone.

\begin{figure*}[t]
 \centering
 \includegraphics[width=0.49\textwidth]{fig6_fullq_positive.pdf}\hfill
 \includegraphics[width=0.49\textwidth]{fig6_fullq_negative.pdf}
 \caption{Complete connected $S_T^{zz}(q,\omega)$ for $L=40$ open chains at
 $J_2/J_1=0.5$ (left) and $-0.5$ (right), with $J_3/J_1=0.5$.
 The physical momentum spans $0\le q\le2\pi$ and energies are in units of $J_1$.
 The triangular time window at $TJ_1=60$ has width $0.09277J_1$;
 production settings are $\Delta tJ_1=0.05$ and $\chi=192$.
 Both panels use the same blue-to-red color coordinate
 $u(s)=\operatorname{asinh}(s/0.02)/\operatorname{asinh}(6/0.02)$ for
 $s=J_1S_T^{zz}\in[0,6]$.
 This uniform nonlinear display enhances weak weight without changing the
 spectral values or their energy resolution. No Gaussian smoothing,
 momentum interpolation, or reflection averaging is used.
 The panels render the calculated
 spectral arrays.}
 \label{fig:fullq}
\end{figure*}

\subsection{Finite-size reconstruction on the ferromagnetic side}
\label{sec:negative}

The static data ask a complementary question: as $\alpha$ approaches zero from
below at fixed $\beta$, does the dominant response retain the $0.60\pi$
momentum favored among the weak-limit antiferromagnetic replicas?
The negative-$\alpha$ scan instead finds a movement toward $q=\pi$.
Figure~\ref{fig:negative-reconstruction} displays the converged $L=80$
samples and their $L=60$ comparison, without interpolating a phase surface.

To test whether this movement locates a common boundary, we require an adjacent
coupling interval to show both a change in $\qpeak$ and a sign change of
$\Delta R$, with $|\Delta R|\ge2\epsilon$ at both endpoints.
Here $\epsilon$ is a deterministic numerical tolerance, built from the
initial-state spread, truncation error, and a $10^{-4}$ floor; it is not a
statistical standard deviation.
Of eight candidate intervals across seven $\beta$ slices, seven retain at
least one signature, but none has both. Only the $\beta=0.05$ momentum
interval reaches the target width $0.025$.
The $\beta=0.10$ slice supplies no supported internal interval at the final
size pair.

The dominant momentum is therefore changing in these finite chains, while the
two boundary diagnostics do not coincide.
This establishes a finite-size reconstruction rather than a separate
thermodynamic phase. A broad crossover, slow convergence of an incommensurate
regime, and residual finite-size effects remain possible interpretations.
An apparent entropy suppression near $(\alpha,\beta)=(-0.10,0.25)$ also fails
an eigensolver control: the fitted coefficient changes from about $0.53$ to
$0.93$ at $L=80$ (Appendix~\ref{sec:krylov}).
It therefore supplies no independent evidence for a new low-entropy phase.

\begin{figure*}[t]
 \centering
 \includegraphics[width=0.96\textwidth]{fig8_finite_size_reconstruction.pdf}
 \caption{Finite-size momentum reconstruction for negative $\alpha$.
 (a) Dominant momentum at the 22 sampled $L=80$ couplings; crosses mark
 disagreement of the momentum classification between $L=60$ and 80.
 Gray segments connect sampled points only.
 (b) Peak-sharpness flow $\Delta R=R_{80}-R_{60}$.
 Bars show $\epsilon=(\epsilon_{80}^2+\epsilon_{60}^2)^{1/2}$, where each
 size-specific tolerance is the larger of the initial-state half-range,
 ten times the maximum final truncation error, and $10^{-4}$.
 These are not statistical confidence intervals.
 (c) Eight candidate intervals tested for a momentum change, a sign change
 satisfying $|\Delta R|\ge2\epsilon$ at both endpoints, width at most $0.025$,
 and the joint momentum/size-flow criterion. No target satisfies the joint
 criterion; this operational test neither establishes nor rules out a
 thermodynamic transition. The symbols
 and decisions are taken from the recorded numerical analysis.}
 \label{fig:negative-reconstruction}
\end{figure*}

\section{Discussion}
\label{sec:discussion}

\subsection{Internal spin structure and momentum-resolved weight}

The projected response in Eq.~\eqref{eq:folded-response} separates two
physical ingredients. The effective chain determines the collective spectrum
at cell momentum $Q$, while the five-site spin distribution determines its
visibility at physical momentum $q$ through $F(q)/5$.
Consequently, a weak response at one replica need not indicate a missing
collective mode. At negative $\alpha$, the near cancellation of the
$q=\pi$ cell amplitude coexists with a strong amplitude at $q=0.6\pi$,
although both momenta fold to $Q=\pi$.
At positive $\alpha$ the same comparison favors $q=\pi$.
The switch of the isolated cell from spin $1/2$ to spin $3/2$ changes this
interference pattern as well as the scale of the inter-cell exchange.

The two-cell calculation tests this statement without a phenomenological
line shape or an intensity fit. At $\alpha=\pm0.5$ and
$\beta/\Delta_{\mathrm{loc}}\le0.01$, the largest relative error of the
41-momentum low-band weight vector is $0.27\%$, and that of the
singlet--triplet gap is $0.56\%$.
These checks validate the absolute projected weights for the specified
finite system in the weak-coupling limit.
For longer chains the same leading factor $F(q)/5$ multiplies the effective
response, but the two-cell line spectrum and its integrated absolute
weights cannot be substituted for those of an extended chain.

This use of the cluster projection complements the trimer and tetramer
studies, which compared numerical spectra with effective descriptions of
propagating excitations~\cite{cheng2022trimer,li2026tetramer}.
Here the quantitative comparison concerns the low-energy intensity associated
with two different half-integer cell spins and their folded momentum
replicas. It does not identify the high-energy bands of
Fig.~\ref{fig:fullq} with unique particles.
Indeed, the two plotted points have
$\beta/\Delta_{\mathrm{loc}}=1.46$ and $3.96$ for positive and negative
$\alpha$, respectively. The two-cell controls at these points already fail
the separation required by the weak-coupling low-band construction.
Their $L=40$ spectra therefore constitute a finite-coupling comparison,
not an extrapolation of the absolute weak-coupling formula.

\subsection{A different local limit near zero intra-cell exchange}

The quartet projection for negative $\alpha$ and the reconstruction toward
$q=\pi$ near $\alpha=0$ are not contradictory: the projection presumes that
the isolated-cell gap is large compared with $\beta$.
At fixed positive $\beta$ this condition fails as $\alpha$ approaches zero,
and a different partition of the chain becomes appropriate.
In the bulk at $\alpha=0$, sites 2--4 of each cell form a spin-$1/2$ trimer,
while the bond joining site 5 to site 1 of the next cell forms a singlet
dimer. Open boundaries additionally leave uncoupled end spins at this
decoupled point.

For $|\alpha|\ll\min(1,\beta)$ the trimer doublets interact through a virtual
triplet of the intervening dimer. Each trimer endpoint projects onto
$2\bm T/3$, and eliminating the dimer triplet yields
\begin{equation}
\Jeff=\frac{2\alpha^2}{9\beta}>0
\label{eq:near-zero}
\end{equation}
in units of $J_1$.
Thus both signs of sufficiently small $\alpha$ lead, at this order, to
antiferromagnetically coupled trimer doublets.
In an open chain, however, each unpaired end spin couples to its adjacent
trimer already at first order, as
$(2\alpha/3)\bm S_{\mathrm{edge}}\!\cdot\!\bm T$.
This sign-dependent $O(\alpha)$ boundary term is larger than the bulk
$O(\alpha^2/\beta)$ scale in the stated limit.
Equation~\eqref{eq:near-zero} is therefore not a formula for the lowest
sector gap of the complete open chain.
Exact diagonalization of the eight-spin trimer--dimer--trimer cluster
supports the coefficient and the $\alpha^2/\beta$ scaling in the
asymptotic limit. The corresponding local-gap estimates at
$(\alpha,\beta)=(-0.10,0.25)$ and $(-0.175,0.25)$ are already
$36.1\%$ and $53.9\%$ below the leading prediction, respectively.
Equation~\eqref{eq:near-zero} therefore explains the limiting mechanism
and its small energy scale, not the finite-coupling location of the
momentum reconstruction.

The static scan compares these competing local descriptions in finite
chains. Its dominant momentum changes from about $0.6\pi$ toward $\pi$,
but the momentum change and the size trend of the peak sharpness do not
select the same coupling intervals.
A change in the strongest physical-momentum replica alone is not a
thermodynamic phase criterion: it can reflect changes of matrix elements
and short-distance correlations as well as changes of long-distance order.
The data establish the finite-size reconstruction but do not determine
whether any part of it becomes singular in the bulk limit.
The small scale in Eq.~\eqref{eq:near-zero} also explains why a low absolute
energy variance need not by itself control the entropy or a sector gap near
this limit; the solver-sensitivity checks must be judged against the relevant
low-energy scale.

\subsection{What a spin probe can distinguish}

There are two separate uses of interference in the present analysis.
Within the cell ground multiplet, $F(q)$ predicts the relative low-energy
weights of folded replicas. Between different isolated-cell multiplets,
reflection parity can forbid a transition exactly.
For the positive-$\alpha$ pentamer, the $q=\pi$ spin operator is even under
reflection, so the lowest odd-parity excited doublet is dark and higher
allowed sectors carry the local inelastic weight.
Coupling the cells does not conserve the reflection of each cell
independently. The isolated-cell dark state is therefore a weak-coupling
selection rule to be followed under mixing, not an exact prohibition of an
entire many-body band at arbitrary $\beta$.

These distinctions suggest direct tests of a weakly coupled pentamer
realization: compare low-energy intensities at momenta sharing the same
$Q$, and compare the lowest local excitation energy with the first allowed
single-spin signal. Quantitative comparison with a material additionally
requires its positions, exchange constants, and scattering operator.
For an SU(2)-invariant singlet the three Cartesian spin responses are equal,
so $S^{zz}$ determines the spin-correlation tensor before the experimental
form factors and polarization factors are applied.
The common finite-time resolution of Fig.~\ref{fig:fullq},
$0.09277J_1$, is broader than the separately calculated finite-chain gaps.
The maps compare resolved intensity distributions; they do not resolve
those gaps from zero or establish the asymptotic widths of spectral features.

\section{Conclusion}
\label{sec:conclusion}

The Heisenberg pentamer chain connects a change of the local composite spin
to a quantitative redistribution of low-energy spectral weight.
In the isolated-cell ranges studied, positive and negative outer intra-cell
exchange select spin-$1/2$ and spin-$3/2$ ground multiplets, respectively.
Projecting the physical spins onto these multiplets fixes both the leading
antiferromagnetic inter-cell exchange and the momentum-dependent interference
factor. The latter favors different physical replicas of the same
effective-chain antiferromagnetic momentum: $q=\pi$ for $\alpha=0.5$ and
$q=0.6\pi$ for $\alpha=-0.5$.
Two-cell exact diagonalization verifies the resulting absolute low-band
weights and excitation scale in the specified weak-coupling regime.
Reflection parity independently explains why the lowest excited doublet of
the positive-$\alpha$ isolated cell is invisible to a $q=\pi$ spin probe.

Finite-chain calculations place these local predictions in a broader
coupling range. Positive-coupling static diagnostics support an effective
antiferromagnetic chain at the tested points, while negative couplings
produce a size-dependent reconstruction of the dominant static momentum.
The alternative trimer--dimer limit near $\alpha=0$ explains why the
isolated-quartet description need not persist there.
Complete connected $L=40$ spectra at $(\alpha,\beta)=(\pm0.5,0.5)$ compare
the responses at one common resolution, outside the controlled
weak-coupling regime. Together, the results distinguish a change in the
visibility of collective spin fluctuations from a demonstrated phase
transition or a unique assignment of high-energy bands.


\begin{acknowledgments}
[Add funding, computing resources, institutional support, and individual
acknowledgments.]
\end{acknowledgments}

\section*{Data and code availability}
The numerical tables underlying the figures, spectral arrays, calculation
and plotting scripts, environment records, and a file-checksum index are
included in the accompanying source-data archive. This archive has not yet
been deposited in a public repository; no public accession identifier is
claimed.

\appendix
\section{Local selection rules and energy scales}
\label{sec:local-details}

Figure~\ref{fig:selection} collects the local projection and selection-rule
data for positive $\alpha$.
At $\alpha=0.14,0.23,0.33$, the endpoint projections give
$\Jeff/\beta=0.440,0.432,0.419$, respectively.
The lowest odd-parity excited doublet carries vanishing $q=\pi$ weight.
Weights are summed over complete final multiplets and averaged over the
initial ground multiplet; the $84\%$--$98\%$ quartet fractions quoted in the
main text refer to the local inelastic response, excluding transitions
within the ground multiplet.

The low-energy effective spin-$1/2$ chain and these local excitations occupy
different scales in the controlled limit.
The effective two-spinon continuum has edges
$\omega_{\mathrm L}(Q)=\pi\Jeff|\sin Q|/2$ and
$\omega_{\mathrm U}(Q)=\pi\Jeff|\sin(Q/2)|$
\cite{descloizeaux1962,caux2006fourspinon}.
Equating its $Q=\pi$ upper edge to the lowest local quartet gap gives formal
values $\beta=0.104,0.176,0.265$ at the three $\alpha$ values.
Their ratios to the cell gap are already $g=1.11,1.14,1.18$.
Such an equality lies outside a small-$g$ expansion and cannot serve as a
controlled prediction of hybridization or of a spectral phase boundary.

\begin{figure*}[t]
 \centering
 \includegraphics[width=0.92\textwidth]{fig5_selection_rules.png}
 \caption{Local projection and selection rules for positive $\alpha$.
 The endpoint coefficients determine the projected exchange, while
 multiplet-resolved weights distinguish bright spin-$3/2$ channels from the
 dark odd-parity doublet at $q=\pi$.
 The connected points in the momentum-weight panel are the five evaluated
 momenta, not a continuous-momentum calculation. No many-body continuum
 boundary or high-energy band assignment is plotted.
 }
 \label{fig:selection}
\end{figure*}

\section{Positive-coupling entropy calibration}
\label{sec:positive-regime}

Entropy offers an independent check of the low-energy interpretation, but the
finite-size slope must be distinguished from a bulk central charge.
At $(\alpha,\beta)=(0.23,0.05)$ we compare the entropy at inter-cell cuts of
$N=8,12,16,20,24$ pentamers ($L=40$--120) with a uniform spin-$1/2$
Heisenberg chain of $N$ effective sites.
The matching is in the number of effective spins, not the physical-site
length. Ratios of the fitted logarithmic slopes are $1.0331$ for the complete
size sequence and $1.0240$ for $N\ge12$; the matched $N=24$ profile gives
$1.0275$ (Fig.~\ref{fig:entanglement}).
The maximum relative entropy change between the tested MPS budgets is
$2.40\times10^{-4}$.
These calibrated values support consistency with the $c=1$ reference.
The raw absolute logarithmic fits fail their fit-quality test and are not
used as precise thermodynamic estimates.

\begin{figure*}[t]
 \centering
 \includegraphics[width=0.94\textwidth]{fig4_entanglement.png}
 \caption{Local-cut entanglement spectra and effective-chain calibration.
 The upper two rows show the lowest relative entanglement levels at five
 inequivalent cuts for $\alpha=0.14$ and $0.33$, both at $\beta=0.5$ and $L=80$.
 The lower comparison is a separate data set at
 $(\alpha,\beta)=(0.23,0.05)$: slopes for pentamer inter-cell cuts are divided
 by those of a uniform spin-$1/2$ chain with the same number of effective
 sites. The three estimators agree with the $c=1$ reference within $3.31\%$.
 The cut spectra are not used as a stand-alone phase classification.
 }
 \label{fig:entanglement}
\end{figure*}

\section{Entropy sensitivity to the local eigensolver}
\label{sec:krylov}
\label{sec:saved-state}

For a general open-chain profile, we fit
\begin{equation}
 S(b)=a_\nu+\frac{\ceff}{6}
 \ln\left[\frac{2L}{\pi}\sin\left(\frac{\pi b}{L}\right)\right],
 \label{eq:entropy}
\end{equation}
with bond-class intercepts $\nu=b\bmod10$ and central windows
$fL\le b\le(1-f)L$, $f=0.10,0.15,0.20,0.25$
\cite{calabrese2004}.
The slope $\ceff$ is descriptive unless finite-size and solver controls
support a central-charge interpretation.

The historical $k=3$ sequence at $(\alpha,\beta)=(-0.10,0.25)$ gives
mean coefficients $0.988,0.695,0.533$ at $L=40,60,80$ for $f=0.20$.
Here $k$ is the local eigensolver's Krylov-subspace dimension, distinct from
the MPS bond dimension $\chi$.
Changing $k=3$ to $8$ at $L=80$ raises the last coefficient to $0.93246$;
the two initial states then differ by only $3.39\times10^{-8}$.
The neighboring $\alpha=-0.175$ control changes from $0.91657$ to $0.95990$.
Figure~\ref{fig:negative-size} keeps these separate fixed-size controls
distinct from the historical size sequence.

At $\alpha=-0.10$, a further $k=8\to12$ control uses both initializations,
64 sweeps, $\chi_{\max}=1200$, and cutoff $10^{-11}$.
The variance per site is $1.9\times10^{-11}$ and the sector gap is
$7.2772\times10^{-5}J_1$.
Across 172 matched observables the changes are bounded by
$7.5\times10^{-12}$ for the gap, $8.4\times10^{-8}$ for the dominant ratio,
and $5.7\times10^{-8}$ for the entropy; $\qpeak=\pi$ is unchanged.
The saved run records and comparison tables have been recovered from the
project archive and checked against their historical checksums.
The anomalously small entropy coefficient does not survive this solver test.
Because the smaller historical sizes still use $k=3$, we do not join them
to the new $L=80$ endpoint as a controlled finite-size extrapolation.

\begin{figure*}[t]
 \centering
 \includegraphics[width=0.96\textwidth]{fig9_entropy_solver_control.pdf}
 \caption{Entropy-coefficient sensitivity at $\beta=0.25$.
 All fits use period-10 bond-class intercepts.
 (a) Historical $k=3$ results at $L=40,60,80$.
 Symbols average the two initial states at $f=0.20$; ranges cover both
 initializations and all four windows.
 (b,c) Separate $L=80$ comparisons of $k=3$ and $8$.
 Symbols show initial-state means at each window and bars the corresponding
 ranges; the two solver series are shifted horizontally by $\mp0.002$ for
 visibility. The dotted line is a reference value, not a bulk central-charge
 determination. Ranges describe deterministic sensitivity, not confidence
 intervals. No $k=12$ values are plotted. The displayed values are taken from the numerical fit tables.}
 \label{fig:negative-size}
\end{figure*}

\section{Spectral resolution and numerical controls}
\label{sec:spectral-control}

The ground-state energies used for the two production spectra are
$E_0/J_1=-12.4925790712$ and $-11.4821383028$, with variance per site
$5.55\times10^{-12}$ and $5.36\times10^{-12}$.
The source norm is checked against the independently contracted equal-time
connected correlations. Complete $L=10$ exact-diagonalization spectra,
processed with the same finite-time estimator, agree with MPS spectra to
relative $L_2$ errors $7.46\times10^{-5}$ and $3.59\times10^{-5}$.
This is a direct small-system validation of the time-evolution and Fourier
normalization, separate from the low-band projection benchmark.

We perform single-variable checks at
$q/\pi=0.2,0.6,1.0$ for both signs of $\alpha$, using the same saved ground
states and $T=60/J_1$.
First, $\chi$ is increased from 192 to 256 at fixed $\Delta t=0.05/J_1$;
then $\Delta t$ is halved at fixed $\chi=256$.
The latter comparison reuses the matching $q=\pi$ record from an earlier
simultaneous refinement of both parameters.
Table~\ref{tab:independent-controls} reports the maximum relative
$L_2$ changes over the three momenta, with the refined spectrum as denominator.
The low-frequency norm is restricted to $0\le\omega\le0.6J_1$.
All 12 comparisons satisfy the criteria fixed before these calculations:
full- and low-frequency changes below $3\%$ and $5\%$, main-peak shifts
below $0.015J_1$, and relative norm drift below $1\%$.
For each of the bands $[0,0.6]$, $[0.6,2.5]$, and $[2.5,4]$ in units of
$J_1$, the weight criterion is
$|W-W_{\rm ref}|\le\max(0.05|W_{\rm ref}|,0.001C_q(0))$.
The largest observed relative band-weight change is $0.376\%$, the largest
norm drift is $0.338\%$, and none of the sampled main peaks changes its
frequency-grid index. These tests constrain the numerical sensitivity of
selected finite-$L$, finite-$T$ spectra; they do not establish long-time,
all-momentum, or thermodynamic convergence.

\begin{table}[!htbp]
 \caption{Independent numerical controls at $L=40$ and $T=60/J_1$.
 Entries are maximum relative $L_2$ changes in percent over
 $q/\pi=0.2,0.6,1.0$. The two columns use $[0,4]$ and $[0,0.6]$ frequency
 intervals in units of $J_1$.}
 \label{tab:independent-controls}
 \begin{ruledtabular}
 \begin{tabular}{rlrr}
 $\alpha$ & Changed parameter & Full & Low\\
 $+0.5$ & $\chi:192\to256$ & 0.0124 & 0.0126\\
 $-0.5$ & $\chi:192\to256$ & 0.0407 & 0.0149\\
 $+0.5$ & $J_1\Delta t:0.05\to0.025$ & 0.8100 & 0.8156\\
 $-0.5$ & $J_1\Delta t:0.05\to0.025$ & 0.5280 & 0.5346\\
 \end{tabular}
 \end{ruledtabular}
\end{table}

The separately computed $L=40$ sector gaps use joint DMRG budgets
$(\chi,\epsilon_{\mathrm{cut}})=(512,10^{-12})$ and $(768,10^{-13})$.
Their changes are of order $10^{-12}J_1$, much smaller than the finite-time
spectral width. Agreement of these two budgets does not make the gaps
thermodynamic quantities or provide a rigorous error estimate.
Similarly, the display step $\Delta\omega=0.005J_1$ does not reduce the
$0.09277J_1$ window width; the heatmaps do not establish sub-linewidth gaps.

For the two heatmaps a single global intensity map is used:
$u(s)=\operatorname{asinh}(s/0.02)/\operatorname{asinh}(6/0.02)$ for
$s=J_1S_T^{zz}\in[0,6]$.
The same map is applied at all $q$ and $\omega$, without energy-dependent
enhancement or rescaling of individual columns.
This is a disclosed color normalization; the accompanying arrays retain
the original linear intensities.
It enhances finite-window sidelobes as well as genuine weak spectral weight,
which is why band identification cannot rest on visual contrast alone.

\subsection{How much weight lies outside the displayed frequency range?}

The finite-time estimator spreads positive-energy lines across zero frequency.
To quantify this effect without modifying the maps, we transform the saved
$C_q(t)$ on a zero-padded discrete Fourier grid with $N_\omega=65536$.
The complete discrete frequency period is
$[-\pi/\Delta t,\pi/\Delta t)$.
The identity $\sum_k S_T^{zz}(q,\omega_k)\Delta\omega=C_q(0)$ is an
algebraic check of this transform, not an independent dynamical sum-rule test
or a proof of infinite-frequency convergence.
An independent contraction of the static connected correlations agrees with
$C_q(0)$ within $5.60\times10^{-11}$ in relative error for nonzero sources.
Direct quadrature on the original display grid reconstructs the saved
spectra with relative error of order $10^{-15}$.

Table~\ref{tab:window-weight} splits the signed spectral integral, normalized
by $C_q(0)$, into negative frequencies, the displayed range, and positive
frequencies above $4J_1$ within the discrete period.
Half the zero-frequency node enters each side; the negative Nyquist node
belongs to the negative side. No negative spectral values are clipped.
Changing $N_\omega$ from 32768 to 65536 changes these fractions by at most
$1.09\times10^{-4}$ over the nonzero-source momentum grid.
This tests numerical quadrature of the same time record, not time resolution.
At $\alpha=-0.5$, the $0.20\pi$ and $0.60\pi$ responses have approximately
$28\%$ and $23\%$ of their weight below zero after window broadening.
These are not physical negative-energy excitations; they demonstrate why
the visible positive-frequency area cannot be identified with the complete
static weight or repaired by stronger color enhancement.

\begin{table}[!htbp]
 \caption{Fractions of the finite-window spectrum, as percentages of $C_q(0)$.
 The intervals refer to the discrete frequency period, with zero split equally
 between the first two columns of weights.}
 \label{tab:window-weight}
 \begin{ruledtabular}
 \begin{tabular}{rrrrr}
 $\alpha$ & $q/\pi$ & $f_-$ & $f_{[0,4]}$ & $f_{>4}$\\
 $+0.5$ & 0.2 & 1.158 & 98.684 & 0.158\\
 $+0.5$ & 0.6 & 0.722 & 99.085 & 0.193\\
 $+0.5$ & 1.0 & 5.683 & 94.168 & 0.149\\
 $-0.5$ & 0.2 & 27.833 & 72.027 & 0.140\\
 $-0.5$ & 0.6 & 22.807 & 77.041 & 0.152\\
 $-0.5$ & 1.0 & 6.798 & 93.022 & 0.180\\
 \end{tabular}
 \end{ruledtabular}
\end{table}


% 先排完附录图表，再开始参考文献。
\clearpage

% 压缩参考文献条目间距。
\setlength{\bibsep}{2.0pt plus 0.3ex}
\bibliography{references}

\end{document}
